% MT005 / answer-zh — canonical editable content.
\textbf{随卷答案选A、B、C、D。} NB先估计先验和类条件分布，再推导后验。对二分类标签\(Y\in\{0,1\}\)，令\(p=P(Y=1\mid x)\)，则\(\operatorname{Var}(Y\mid x)=p(1-p)\)在\(p=1/2\)时最大；这对应通常的等错分代价决策边界。非高斯的NB可以产生非线性边界，NB和逻辑回归都可以用最大似然估计训练。E错误，因为平滑可以对NB起到正则化作用。A描述的是最终得到的后验，并不表示NB直接采用判别式训练。
